
\chapter{NumPy package}
%%%%%%%%%%%%%%%%%%%%%%%%
% $Autor: MansourAhmadyPhoulady, RanjeetsingTawar, HemanthPrabhakaran$
% $Pfad: Car_Licence_Plate_Identification_with_a_JetsonNano$
% $Version: 4.0.0 $
%
% !TeX encoding = utf8
%
%%%%%%%%%%%%%%%%%%%%%%%%

\section{Introduction}

NumPy, short for Numerical Python, is a widely-used open-source Python library that serves as a cornerstone in various scientific and engineering fields. Renowned for its robust capabilities in handling numerical data, it has become the universal standard for numerical computing in Python and forms the foundation of the scientific Python and PyData ecosystems. NumPy is utilized by a diverse audience, ranging from novice programmers to seasoned researchers engaged in cutting-edge scientific and industrial research. Moreover, its API is a critical component of many prominent Python libraries, including Pandas, SciPy, Matplotlib, scikit-learn, scikit-image, and other data science tools.\cite{NumPy:2024}

At the heart of NumPy lies its support for multidimensional arrays and matrix data structures, encapsulated in its efficient ndarray object. This object allows users to perform a vast range of mathematical operations with ease. NumPy empowers Python with powerful data structures designed for efficient computations involving arrays and matrices, complemented by an extensive library of high-level mathematical functions. These features not only enhance computational efficiency but also simplify the implementation of complex mathematical models and algorithms.

In addition to its core functionality, NumPy plays a pivotal role in fostering collaboration and innovation within the data science and scientific computing communities. Its interoperability with other Python libraries ensures seamless workflows, making it indispensable for tasks like data preprocessing, numerical simulations, and machine learning. By bridging the gap between high-level programming and optimized performance, NumPy enables researchers and developers to focus on solving real-world problems effectively. \cite{NumPy:2024}

\section{Description}

NumPy is an essential library for scientific computing in Python. It provides a powerful framework that includes a multidimensional array object, derived objects such as masked arrays and matrices, and a wide range of functions for fast operations on arrays. These operations cover various mathematical, logical, and statistical tasks, shape manipulation, sorting, selection, I/O, discrete Fourier transforms, basic linear algebra, random simulations, and more \cite{NumPy:2024}.

The core of the NumPy library is the ndarray object, which supports n-dimensional arrays of homogeneous data types. Many of its operations are executed in compiled code, which significantly improves performance. NumPy arrays differ from standard Python sequences in several key ways:

\begin{itemize}
	\item \textbf{Fixed Size:} Unlike Python lists that can dynamically grow, NumPy arrays have a fixed size upon creation. Changing the size of an ndarray will create a new array, deleting the original one \cite{NumPy:2024}.
	\item \textbf{Homogeneous Data Types:} All elements in a NumPy array must be of the same data type, which ensures that they occupy the same amount of memory. However, arrays of Python objects (including NumPy objects) can store elements of varying sizes.
	\item \textbf{Efficiency:} NumPy arrays are optimized for advanced mathematical and other operations on large datasets. These operations are typically faster and require less code compared to using Python’s built-in sequences.
\end{itemize}

Many scientific and mathematical Python packages rely on NumPy arrays, converting Python sequence inputs into NumPy arrays before processing and often outputting NumPy arrays as well. Consequently, to effectively use a significant portion of modern scientific and mathematical Python software, it is crucial to understand how to work with NumPy arrays, beyond just using Python’s built-in sequence types.

\subsection{Why is NumPy Fast?}

Vectorization refers to the absence of explicit loops, indexing, or other manual operations in the code. Although these actions occur in the background, they are executed efficiently through optimized, pre-compiled C code. The benefits of vectorized code include:

\begin{itemize}
	\item \textbf{Conciseness and Readability:} Vectorized code is more compact and easier to understand.
	\item \textbf{Reduced Bugs:} With fewer lines of code, there are typically fewer chances for errors.
	\item \textbf{Mathematical Precision:} The code often mirrors standard mathematical notation, making it simpler to implement mathematical operations correctly.
	\item \textbf{More Pythonic Code:} Vectorization leads to cleaner, more efficient code compared to code cluttered with loops that are harder to read and understand.
\end{itemize}

Broadcasting is a concept in NumPy that refers to the implicit element-by-element operation behavior. It is used not only for arithmetic but also for logical, bitwise, and functional operations. In NumPy, many operations operate element-wise by default, meaning that even if arrays have different shapes, as long as they can be expanded to match, the operation can be performed correctly. This flexibility allows for the use of scalars, arrays, or multidimensional arrays in operations without needing explicit loops \cite{NumPy:2024}.

\subsection{Key Features of NumPy}

\begin{itemize}
	\item \textbf{N-dimensional array object (ndarray):} NumPy provides the ndarray object, which is a versatile and efficient array structure designed to store and manipulate large datasets across multiple dimensions.
	\item \textbf{Vectorized operations:} NumPy supports efficient computation by enabling operations on entire arrays simultaneously, eliminating the need for manual iteration over individual elements.
	\item \textbf{Broadcasting:} With broadcasting, NumPy allows arrays of different shapes to be combined in arithmetic operations, making the code more concise and easier to understand.
	\item \textbf{Element-wise operations:} NumPy offers a wide range of mathematical functions that perform element-wise operations, such as arithmetic, trigonometric, and logarithmic functions, on arrays.
	\item \textbf{Array indexing and slicing:} NumPy provides advanced indexing and slicing techniques, enabling efficient access and manipulation of array elements, which enhances data processing.
	\item \textbf{Linear algebra operations:} NumPy includes built-in functions for essential linear algebra tasks like matrix multiplication, matrix decomposition, eigenvalue computation, and solving linear equations.
	\item \textbf{Random number generation:} NumPy offers tools for generating random numbers from different distributions, as well as functions for shuffling arrays and generating random samples \cite{NumPy:2024}.
	\item \textbf{Integration with other libraries:} NumPy integrates smoothly with other Python libraries used in scientific computing, such as SciPy, matplotlib, and pandas, facilitating complex computational workflows.
	\item \textbf{Efficient memory management:} NumPy's array structure is implemented in C, which ensures efficient memory handling and optimizations, making it well-suited for managing large datasets.
	\item \textbf{Open-source and community-driven:} As an open-source project, NumPy benefits from a large and active community, which ensures its ongoing development and improvement with regular updates.
\end{itemize}

\section{Installation}

To install NumPy, the only requirement is Python. If you haven't installed Python yet and are looking for an easy way to get started, we recommend using the Anaconda Distribution, which includes Python, NumPy, and various other popular packages for scientific computing and data science.

NumPy can be installed using methods like \texttt{conda}, \texttt{pip}, package managers on macOS and Linux, or by compiling from the source. For more comprehensive installation instructions, refer to below Python and NumPy installation guide.


\subsection{Python package management}
Package management can be a complex task, which has led to the development of various tools. For general Python development, there are numerous tools that complement pip, and for high-performance computing (HPC), Spack is a useful option. However, for most NumPy users, \textbf{Pip \& conda} are the two most commonly used tools.

\bigskip

Package management can be a complex task, which has led to the development of various tools. For general Python development, there are numerous tools that complement pip, and for high-performance computing (HPC), Spack is a useful option. However, for most NumPy users, conda and pip are the two most commonly used tools.

\bigskip

One major difference is that conda is cross-language and can install Python itself, while pip is specific to a particular Python installation and can only install packages for that Python. This means conda can also handle non-Python libraries and tools, such as compilers, CUDA, or HDF5, while pip cannot.

\bigskip

Another distinction is that pip installs packages from the Python Packaging Index (PyPI), which is the largest package repository, whereas conda installs from its own channels, like "defaults" or "conda-forge." While PyPI has a broader selection, popular packages are also available through conda.

\bigskip

Lastly, conda offers an integrated solution for managing packages, dependencies, and environments, whereas with pip, you may need additional tools to manage environments or handle complex dependencies.

\subsubsection{Installing with Anaconda}

If you use conda, you can install NumPy from the defaults or conda-forge channels:

\begin{code}[h!]
	\lstinputlisting[language=Python, numbers=none, linerange={23-29}]{Code/NumPy/NumPy.py}    
	\caption{Installing NumPy using terminal}
	
\end{code}	

\subsubsection{Installing using terminal or command prompt (PIP)}

If you use pip, you can install NumPy with:	

\begin{code}[h!]
	\lstinputlisting[language=Python, numbers=none, linerange={31-32}]{Code/NumPy/NumPy.py}    
	\caption{Installing NumPy using PIP}
\end{code}	


\subsection{Required dependencies}

To install NumPy, the following dependencies are required:

\begin{enumerate}
	\item \textbf{Python:} Since NumPy is a Python library, you must have Python installed on your system. It supports Python versions 3.7 and newer.
	
	\item \textbf{Compiler:} When installing NumPy with pip, a compiler is needed on your system. On Windows, this typically involves installing Microsoft Visual C++ Build Tools or MinGW. On Linux, you may need to install a package like build-essential.
	
	\item \textbf{NumPy:} Finally, you will need to install NumPy itself. This can be done easily using pip, the Python package manager, by running the command \texttt{pip install numpy}.
	
\end{enumerate}

\section{Example - Manual}

\subsection{User Manual for the Example Python File in PyCharm}

\subsection*{Installation and Setup}

\begin{enumerate}
	
	\item \textbf{Download and install Python:}
	
	Visit the official Python website \url{https://www.python.org} and download the most recent version of Python for your operating system. Follow the given instructions to complete the installation process.
	
	\item \textbf{Install PyCharm: }
	
	Visit the JetBrains website \url{https://www.jetbrains.com/pycharm} and download the free PyCharm Community Edition. Follow the installation guide specific to your operating system to set it up.
	
\end{enumerate}

\subsection*{Opening the Python File in PyCharm}

\begin{enumerate}
	
	\item \textbf{Launch PyCharm:} Open the PyCharm application from your desktop or the applications menu.
	
	\item \textbf{Create a new project:} Click on \texttt{"Create New Project"} or go to \texttt{"File"} $\rightarrow$ \texttt{"New Project"}. Choose a appropriate location for your project and provide a name.
		
	\item \textbf{Open the Python file:} After creating the project, navigate to the project directory in the PyCharm project view. Right-click on the desired folder and select \texttt{"New"} $\rightarrow$ \texttt{"Python File"}. Provide a name for the file and click \texttt{"OK"}.
	
	\item \textbf{Copy your Python code:} Open your Python file (with .py extension) in any text editor and copy the contents.
	
	\item \textbf{Paste the code:} Paste the copied code into the newly created Python file in PyCharm.
	
\end{enumerate}

\subsection*{Working with the Python File}

\begin{enumerate}
	
	\item \textbf{Running the script:} To run the Python script, right-click anywhere within the Python file and select \texttt{"Run"} $\rightarrow$ \texttt{"Run ExampleManual.py"}. Alternatively, you can use the keyboard shortcut \texttt{"Shift + F10"}. The script will execute, and the output will be displayed in the PyCharm console.
	
	\item \textbf{Debugging the script:} To debug the script and set breakpoints, click on the left margin of the Python file next to the line where you want to place the breakpoint. A red dot will appear, indicating the breakpoint. Click on the \texttt{Debug} button or use the keyboard shortcut \texttt{"Shift + F9"} to start debugging the script.
	
	\item \textbf{Interacting with the script:} If your script prompts for user input, you can provide the necessary input directly in the PyCharm console. The console enables interaction with the script as it runs.
	
\end{enumerate}

\subsection*{Modifying the Python File}

\begin{enumerate}
	
	\item \textbf{Editing the code:} To modify the Python code, find the section you want to change and make the necessary edits.
	
	\item \textbf{Saving the changes:} PyCharm automatically saves your work as you go. However, you can also save the file manually by selecting \texttt{"File"} $\rightarrow$ \texttt{"Save"} or using the keyboard shortcut \texttt{"Ctrl + S"}.
	
\end{enumerate}

\subsection*{Further Assistance and Resources}

\begin{itemize}
	
	\item \textbf{PyCharm Documentation:} PyCharm offers comprehensive documentation to help you understand its features and functionality. You can access it online at \texttt{https://www.jetbrains.com/help/pycharm}.
	
	\item \textbf{Python Documentation:} The official Python documentation provides detailed information about the Python language, libraries, and best practices. It is available at \url{https://docs.python.org}.
	
	\item \textbf{Online Python Communities:} Joining online communities like Stack Overflow \url{https://stackoverflow.com} or the Python subreddit \url{https://www.reddit.com/r/Python} can provide valuable insights and assistance from experienced Python developers.
	
\end{itemize}

\subsection{NumPy Verison Check}
You can check the version of Numpy by using the following command:

\begin{code}[h!]
	\lstinputlisting[language=Python, numbers=none, linerange={37-39}]{Code/NumPy/NumPy.py}    
	\caption{Version Check}
\end{code}	

\subsection{Example Files}
You can save and load numpy array into a file with numpy functions such as follows: “numpy.save” and “numpy.load”.

\begin{code}[h!]
	\lstinputlisting[language=Python, numbers=none, linerange={41-44}]{Code/NumPy/NumPy.py}    
	\caption{NumPy-Example Files}
\end{code}		

\subsection{How to import}
The code begins by importing the required libraries, numpy (\texttt{np}).

\begin{code}[h!]
	\lstinputlisting[language=Python, numbers=none, linerange={34-35}]{Code/NumPy/NumPy.py}    
	
	\caption{How to Import the package}
	
\end{code}

\subsection{Creating Arrays}

Creating arrays is the foundation of NumPy. You can create arrays using various methods like \texttt{np.array()}, \texttt{np.zeros()}, \texttt{np.ones()}, \texttt{np.arange()}, and \texttt{np.linspace()}.

\begin{code}[h!]
	\lstinputlisting[language=Python, numbers=none, linerange={46-60}]{Code/NumPy/NumPy.py}    
	
	\caption{Creating Arrays}
	
\end{code}

\subsection{Array Operations}

NumPy allows you to perform various mathematical operations on arrays, including arithmetic operations, trigonometric functions, exponential functions, and more.

\begin{code}[h!]
	\lstinputlisting[language=Python, numbers=none, linerange={62-76}]{Code/NumPy/NumPy.py}    
	
	\caption{Array Operations}
	
\end{code}

\subsection{Array Indexing and Slicing:}

NumPy provides powerful indexing and slicing capabilities to access and manipulate elements within arrays.

\begin{code}[h!]
	\lstinputlisting[language=Python, numbers=none, linerange={78-89}]{Code/NumPy/NumPy.py}    
	
	\caption{Array Indexing and Slicing}
	
\end{code}

\subsection{Array Shape Manipulation:}

NumPy provides functions to change the shape and dimensions of arrays.

\begin{code}[h!]
	\lstinputlisting[language=Python, numbers=none, linerange={91-101}]{Code/NumPy/NumPy.py}    
	
	\caption{Array Shape Manipulation}
	
\end{code}

\subsection{Linear Algebra Operations:}

NumPy offers a comprehensive set of functions for linear algebra operations.

\begin{code}[h!]
	\lstinputlisting[language=Python, numbers=none, linerange={103-118}]{Code/NumPy/NumPy.py}    
	
	\caption{Linear Algebra Operations}
	
\end{code}


\subsection{Importing and exporting data}

In NumPy, you can import and export data from various file formats such as text files, CSV files, and binary files.Depending on the data format and requirements, you may need to customize the import/export functions accordingly. Here are examples of importing and exporting data using NumPy:

\begin{code}[h!]
	\lstinputlisting[language=Python, numbers=none, linerange={121-140}]{Code/NumPy/NumPy.py}    
	
	\caption{Importing and exporting data in NumPy}
	
\end{code}



\subsection{NumPy Error Handling}

NumPy, like many other libraries in Python, provides error handling mechanisms to manage exceptions and errors that may occur during numerical computations and array operations. These techniques are crucial for diagnosing and addressing issues that arise during data manipulation and analysis. Below are some common error handling techniques in NumPy:

\begin{itemize}
	\item \textbf{Try-Except Blocks:} Use try-except blocks to catch and handle specific exceptions during NumPy operations, ensuring smooth execution of code even in the presence of errors.
	\item \textbf{Error Reporting:} NumPy provides descriptive error messages and tracebacks to pinpoint the location and nature of errors, aiding in effective debugging and troubleshooting.
	\item \textbf{Error Handling Functions:} Utilize NumPy functions like \texttt{np.seterr()} to control error handling modes, such as raising exceptions, displaying warnings, or ignoring errors altogether.
	\item \textbf{Error Handling with Arrays:} Employ array bounds checking and special value detection functions like \texttt{np.isnan()} and \texttt{np.isinf()} to prevent errors and handle exceptional conditions when working with NumPy arrays.
\end{itemize}


\section{Further Reading}

\begin{enumerate}
	\item \textbf{NumPy User Guide:} \\
	The official NumPy user guide provides comprehensive documentation on NumPy's functionality, including array creation, manipulation, mathematical operations, and more. 
	\href{https://numpy.org/doc/stable/user/index.html}{NumPy User Guide}
	
	\item \textbf{Python Data Science Handbook:} \\
	This book by Jake VanderPlas covers essential tools and techniques for data science in Python, with a significant focus on NumPy. It provides practical examples and explanations of NumPy's usage in data manipulation, analysis, and visualization. 
	\href{https://jakevdp.github.io/PythonDataScienceHandbook/}{Python Data Science Handbook}
	
	\item \textbf{NumPy Official Documentation:} \\
	The official documentation for NumPy is an invaluable resource for learning about NumPy's functions, methods, and usage patterns. It provides detailed explanations and examples for each aspect of the library \cite{NumPy:2024}. 
	\href{https://numpy.org/doc/stable/}{NumPy Official Documentation}
	
	\item \textbf{Numerical Python:} \\
	This book by Robert Johansson offers a comprehensive introduction to NumPy and its applications in scientific computing. It covers fundamental concepts, advanced techniques, and practical examples to help you master NumPy. 
	\href{https://www.apress.com/gp/book/9781484242452}{Numerical Python}
	
	\item \textbf{NumPy Cookbook:} \\
	The NumPy Cookbook by Ivan Idris and Eli Bressert provides a collection of recipes for common tasks and challenges encountered in numerical computing with NumPy. It offers practical solutions and insights for optimizing performance and solving real-world problems. 
	\href{https://www.packtpub.com/product/numpy-1-5-cookbook/9781849515306}{NumPy Cookbook}
\end{enumerate}